\documentclass[11pt]{article}
\usepackage[a4paper,margin=27mm]{geometry}
\usepackage{amsmath,amssymb,amsthm,hyperref}
\newtheorem{theorem}{Theorem}
\title{An exact support bound for uniform triple minima}
\author{Research proof; independently reviewed}
\date{30 September 2026}
\begin{document}\maketitle
Let a probability distribution on permutations of $[n]$ have $q$ support
points, all of positive weight. Require only that each member of every
three-element subset is its earliest element with probability $1/3$.
Pairwise rank balance is not assumed.
\begin{theorem}
For $n\ge4$, one has $q\ge n-1$. Arbitrary positive weights and repeated
permutations are allowed (repetitions can be merged).
The bound is attained at $n=4$ by the three ordered lists
$1203$, $3201$, and $0213$, with equal weights.
\end{theorem}
\begin{proof}
Fix a pivot $i$, and for $j\ne i$ let $v_j$ be the binary indicator of
$i$ preceding $j$, viewed in the $q$-dimensional weighted Euclidean
space of support points. Write $p_{ij}=\Pr(i<j)$.
The Gram matrix on the $m=n-1$ columns is
\[
 G_i=D_i+\frac13\mathbf1\mathbf1^T,\qquad
 (D_i)_{jj}=p_{ij}-\frac13\ge0.
\]
The off-diagonal entries follow from triple-minimum balance. The
inequality on the diagonal follows because being earliest in a triple
containing $i,j$ implies $i<j$.

Suppose $q\le n-2$, so $m>q$. If every diagonal entry of $D_i$ were
positive, then $G_i$ would be positive definite of rank $m$, impossible
in a $q$-dimensional space. Thus some $a\ne i$ has $p_{ia}=1/3$.
For every $j\notin\{i,a\}$,
\[
 \Pr(i<a,\ i>j)=\Pr(i<a)-\Pr(i<a,\ i<j)=0.
\]
Consequently the event $i<a$ implies that $i$ is globally first, up to
a null event. In particular $\Pr(i\text{ globally first})\ge1/3$.
This holds for every pivot $i$. Summing contradicts $n\ge4$.
The displayed example is checked by restricting to its four triples.
\end{proof}

\paragraph{Attribution and scope.}
The condition is called weakly three-restricted min-wise independence.
The argument uses the classical diagonal-plus-constant Gram method.
A 2011 abstract by M.~Vsemirnov states the more general bound
$q\ge(n-k+2)/(k-1)$ for weakly $k$-restricted families; see
\href{https://www.pdmi.ras.ru/EIMI/2011/RuFiDiM/booklet.pdf}{the primary conference booklet},
p.~40. Priority of the sharper special-case application above has not
been established. It is a bound for permutation distributions, not for
general dice. For the concatenated two-face palindrome construction it
implies at least $2n-2$ faces per die whenever triple fairness is exact.
\end{document}
